\chapter{Worked Example: Franka Emika Panda (7-DOF)}
\UseThesisAcronym{DOF}
\label{app:panda}

This appendix contains three parts:
\textbf{Part A} traces one complete 1\,ms control cycle numerically for
a concrete Panda configuration.
\textbf{Part B} provides the complete C++ implementation of Experiment~1
(Cartesian position and orientation tracking) with simulation results.
\textbf{Part C} provides the complete C++ implementation of Experiment~2
(virtual wall constraint) with simulation results.

The numerical worked example in Part~A covers every quantity required for
one 1\,ms cycle — position vector, Jacobian, position-and-orientation error, impedance force,
and joint torques.

% ─────────────────────────────────────────────────────────────────
\section{Robot Parameters and Configuration}
\label{app:panda_params}
% ─────────────────────────────────────────────────────────────────

\subsection{Denavit--Hartenberg Parameters}

The Panda's kinematic chain consists of seven revolute joints.
The Denavit--Hartenberg (\abbr{DH}) parameters published by Franka Emika \cite{FrankaFCI2023}
are reproduced in \cref{tab:panda_dh}.

\begin{table}[htb]
\centering
\caption{DH parameters of the Franka Emika Panda.
         The flange offset (joint\,7 to tool centre) is included in joint\,7.}
\label{tab:panda_dh}
\renewcommand{\arraystretch}{1.3}
\begin{tabular}{@{}ccccc@{}}
\toprule
Joint $i$ & $a_i$ [m] & $d_i$ [m] & $\alpha_i$ [rad] & $\theta_i$ [rad] \\
\midrule
1 & $0$        & $0.333$  & $0$           & $q_1$ \\
2 & $0$        & $0$      & $-\pi/2$      & $q_2$ \\
3 & $0$        & $0.316$  & $+\pi/2$      & $q_3$ \\
4 & $+0.0825$  & $0$      & $+\pi/2$      & $q_4$ \\
5 & $-0.0825$  & $0.384$  & $-\pi/2$      & $q_5$ \\
6 & $0$        & $0$      & $+\pi/2$      & $q_6$ \\
7 & $+0.088$   & $0.107$  & $+\pi/2$      & $q_7$ \\
\bottomrule
\end{tabular}
\end{table}

\subsection{Home Configuration}

The home configuration used throughout this thesis is:

\begin{equation}
\label{eq:q0}
\vq_0 =
\begin{bmatrix}
0 \\ -\pi/4 \\ 0 \\ -3\pi/4 \\ 0 \\ \pi/2 \\ \pi/4
\end{bmatrix}
\approx
\begin{bmatrix}
0^\circ \\ -45^\circ \\ 0^\circ \\ -135^\circ \\ 0^\circ \\ 90^\circ \\ 45^\circ
\end{bmatrix}
\end{equation}

This configuration places the end-effector approximately 0.47\,m in front of
the base with the gripper pointing downward.
It avoids the zero-angle singularity (joints 1, 3, 5 collinear) and provides
good Jacobian conditioning across the workspace.

% ─────────────────────────────────────────────────────────────────
\section{Step 1 --- Forward Kinematics: End-Effector Position}
\label{app:panda_fk}
% ─────────────────────────────────────────────────────────────────

Applying the DH transformation chain of \cref{eq:forward_kinematics_chain} with $\vq = \vq_0$:

\begin{equation}
\mT_0^7(\vq_0) =
\mT_0^1(q_1)\;\mT_1^2(q_2)\;\mT_2^3(q_3)\;\mT_3^4(q_4)\;
\mT_4^5(q_5)\;\mT_5^6(q_6)\;\mT_6^7(q_7)
\end{equation}

The computation is performed numerically by evaluating each $4\times4$ DH
matrix and multiplying left to right.
The resulting end-effector position vector (the controlled variable
from \cref{eq:end_effector_position}) is extracted from the upper-right column of
$\mT_0^7$:

\begin{equation}
\label{eq:pe_panda}
p_{EE}(\vq_0) =
\begin{bmatrix}
0.1164 \\ 0.5452 \\ 0.3952
\end{bmatrix}
\;\mathrm{m}
\end{equation}

Here:

\begin{itemize}
  \item $p_x = 0.1164\,\mathrm{m}$: position along the base-frame $x$-axis.
  \item $p_y = 0.5452\,\mathrm{m}$: position along the base-frame $y$-axis.
  \item $p_z = 0.3952\,\mathrm{m}$: height above the mounting surface.
\end{itemize}

\subsection{Joint-Axis Points and Directions}

\cref{tab:panda_fk} lists, for each joint \(i\), a point
\(p_{i-1}\) on its axis and its rotation axis
\(z_{i-1}\), expressed in the base frame. They are extracted from
\(\mT_0^{i-1}(\vq_0)\) for \(i=1,\ldots,7\), including
\(\mT_0^0=I\) for joint~1, and are used to construct the Jacobian columns.

\begin{table}[htb]
\centering
\caption{Axis points \(p_{i-1}\) and directions
         \(z_{i-1}\) for joints \(i=1,\ldots,7\) at
         \(\vq_0\), using the standard-DH indexing.}
\label{tab:panda_fk}
\renewcommand{\arraystretch}{1.3}
\begin{tabular}{@{}cp{3.0cm}p{1.8cm}p{6.5cm}@{}}
\toprule
Joint \(i\) & $p_{i-1,x}$ [m] & $p_{i-1,y}$ [m] & $p_{i-1,z}$ [m] \\
\midrule
1 & $0.0000$ & $0.0000$ & $0.0000$ \\
2 & $0.0000$ & $0.0000$ & $0.3330$ \\
3 & $0.0000$ & $0.0000$ & $0.3330$ \\
4 & $0.2234$ & $0.2234$ & $0.3330$ \\
5 & $0.1409$ & $0.2234$ & $0.3330$ \\
6 & $0.2234$ & $0.6074$ & $0.3330$ \\
7 & $0.2234$ & $0.6074$ & $0.3330$ \\
\midrule
Joint \(i\) & $z_{i-1,x}$ & $z_{i-1,y}$ & $z_{i-1,z}$ \\
\midrule
1 & $\phantom{-}0.0000$ & $\phantom{-}0.0000$ & $+1.0000$ \\
2 & $\phantom{-}0.0000$ & $\phantom{-}0.0000$ & $+1.0000$ \\
3 & $+0.7071$ & $+0.7071$ & $\phantom{-}0.0000$ \\
4 & $\phantom{-}0.0000$ & $\phantom{-}0.0000$ & $+1.0000$ \\
5 & $\phantom{-}0.0000$ & $+1.0000$ & $\phantom{-}0.0000$ \\
6 & $\phantom{-}0.0000$ & $\phantom{-}0.0000$ & $+1.0000$ \\
7 & $-1.0000$ & $\phantom{-}0.0000$ & $\phantom{-}0.0000$ \\
\bottomrule
\end{tabular}
\end{table}

% ─────────────────────────────────────────────────────────────────
\section{Step 2 --- Jacobian Matrix \texorpdfstring{$J \in \mathbb{R}^{6\times7}$}{J in R(6x7)}}
\label{app:panda_jacobian}
% ─────────────────────────────────────────────────────────────────

Applying \cref{eq:jacobian_column} for each joint using the values from
\cref{tab:panda_fk} and the end-effector position \cref{eq:pe_panda}:

\begin{equation}
\mJ_i =
\begin{bmatrix}
z_{i-1} \times
(p_{\mathrm{EE}} - p_{i-1}) \\[3pt]
z_{i-1}
\end{bmatrix}
\end{equation}

The numerical result is the $6\times7$ geometric Jacobian at $\vq_0$:

\begin{equation}
\label{eq:J_panda}
\mJ(\vq_0) =
\begin{bmatrix}
-0.5452 & -0.5452 &  0.0440 & -0.3218 &  0.0622 &  0.0622 &  0.0000 \\[3pt]
 0.1164 &  0.1164 & -0.0440 & -0.1070 &  0.0000 & -0.1070 &  0.0622 \\[3pt]
 0.0000 &  0.0000 &  0.3032 &  0.0000 &  0.0245 &  0.0000 &  0.0622 \\[3pt]
 0.0000 &  0.0000 &  0.7071 &  0.0000 &  0.0000 &  0.0000 & -1.0000 \\[3pt]
 0.0000 &  0.0000 &  0.7071 &  0.0000 &  1.0000 &  0.0000 &  0.0000 \\[3pt]
 1.0000 &  1.0000 &  0.0000 &  1.0000 &  0.0000 &  1.0000 &  0.0000
\end{bmatrix}
\end{equation}

The rows correspond to $[v_x,\;v_y,\;v_z,\;\omega_x,\;\omega_y,\;\omega_z]$
and the columns to joints 1 through 7.

\begin{remark}
Columns 1 and 2 of $\mJ(\vq_0)$ are identical in this numerical model because
\(z_0=z_1\) at the home configuration and \(p_1-p_0\) is parallel to that
common axis. Their different axis points therefore differ only by a component
that vanishes in the cross product.
This is characteristic of the Panda's geometry at $\vq_0$; it changes with the
robot's configuration.
\end{remark}

% ─────────────────────────────────────────────────────────────────
\section{Step 3 --- Pseudo-Inverse and Null-Space Projector}
\label{app:panda_pinv}
% ─────────────────────────────────────────────────────────────────

The pseudo-inverse $\mJ^+$ is \emph{not} required for the primary torque
computation.
It is computed here only to form the null-space projector in this appendix.

\subsection{\texorpdfstring{Computing $J^+$}{Computing J+}}

For the full-rank Jacobian used in this numerical example:

\begin{equation}
J^+ = \mJ^\top\bigl(\mJ\mJ^\top\bigr)^{-1} \in \R{7\times6}
\end{equation}

The product $\mJ\mJ^\top \in \R{6\times6}$ is formed first, then inverted.
At the home configuration the result is:

\begin{equation}
\label{eq:Jp_panda}
J^+(\vq_0) \approx
\begin{bmatrix}
 0.0000 &  2.2377 &  0.0000 &  0.1392 &  0.0000 &  0.2394 \\[2pt]
 0.0000 &  2.2377 &  0.0000 &  0.1392 &  0.0000 &  0.2394 \\[2pt]
 0.0000 &  0.0000 &  3.0315 &  0.1886 & -0.0743 &  0.0000 \\[2pt]
-2.6042 & -7.0795 &  0.0000 & -0.4405 &  0.1620 & -0.5955 \\[2pt]
 0.0000 &  0.0000 & -2.1436 & -0.1334 &  1.0525 &  0.0000 \\[2pt]
 2.6042 &  2.6042 &  0.0000 &  0.1620 & -0.1620 &  1.1166 \\[2pt]
 0.0000 &  0.0000 &  2.1436 & -0.8666 & -0.0525 &  0.0000
\end{bmatrix}
\end{equation}

\paragraph{Verification.}
The defining property $\mJ\,J^+ = \mI_6$ can be confirmed:

\begin{equation}
\mJ\,J^+ = \mI_6 \quad \checkmark
\end{equation}

\subsection{Null-Space Projector}

The null-space projector is defined as:

\begin{equation}
\label{eq:nullspace}
N_q = \mI_7 - J^+\mJ \in \R{7\times7}
\end{equation}

The diagonal entries of $N_q$ at $\vq_0$ are:

\begin{equation}
\mathrm{diag}(N_q) =
[0.5,\;\; 0.5,\;\; 0.0,\;\; 0.0,\;\; 0.0,\;\; 0.0,\;\; 0.0]
\end{equation}

Here:

\begin{itemize}
  \item Joints\,1 and\,2 have the largest null-space contribution (0.5 each).
        This makes sense: at $\vq_0$ their columns in $\mJ$ are nearly
        identical (see the remark above), so a motion along their
        \emph{difference} direction produces no end-effector velocity.
  \item Joints\,3–7 have zero diagonal entries in $N_q$ at this
        configuration, meaning their individual motions are fully
        task-relevant and cannot be independently redirected without
        affecting the end-effector.
\end{itemize}

The implemented controller uses a projected null-space torque as described in
\cref{sec:nullspace_conditioning_implementation}. In this numerical example,
the projector $N_q$ is computed only to illustrate the structure of the
redundancy.
Physically, the non-zero entries in $\mathrm{diag}(N_q)$ at joints
1 and 2 mean that a motion along the \emph{difference} of those two joint
angles moves no Cartesian end-effector coordinate --- it is the one degree
of freedom that is ``invisible'' to the task.

% ─────────────────────────────────────────────────────────────────
\section{Step 4 --- Impedance Parameter Matrices}
\label{app:panda_matrices}
% ─────────────────────────────────────────────────────────────────

For the Panda, the controlled variable is the full end-effector position and orientation
$T_{EE}\in SE(3)$ with translational error
$e_p\in\mathbb{R}^3$ and rotational error
$e_R\in\mathbb{R}^3$ (see \cref{subsec:cartesian_pose_error}).
The impedance matrices are $3\times3$ symmetric positive definite matrices as
defined in
\cref{sec:stiffness_damping_representation}.

\subsection{\texorpdfstring{Virtual Inertia Matrix $M$}{Virtual Inertia Matrix M}}

\begin{equation}
\label{eq:M_panda}
M = \mI_6 =
\mathrm{diag}(1,\,1,\,1,\,1,\,1,\,1)
\end{equation}

\begin{itemize}
  \item $M_{11} = M_{22} = M_{33} = 1\,\mathrm{kg}$: virtual translational mass.
  \item $M_{44} = M_{55} = M_{66} = 1\,\mathrm{kg\,m^2}$: virtual rotational inertia.
  \item Setting $\mM = \mI_6$ decouples all axes in the simplified
        spring-damper model. The real link inertia and Coriolis effects are handled separately by
        $\bm{\tau}_g$ and $\bm{\tau}_c$ (see \cref{eq:commanded_torque}).
\end{itemize}

\subsection{\texorpdfstring{Stiffness Matrices $K_p$ and $K_R$}{Stiffness Matrices Kp and KR}}

For the axis-aligned virtual wall example, the stiffness values in the wall
frame are chosen as:
$k_{\mathrm{free}}=50\,\mathrm{N/m}$,
$k_{\mathrm{stiff}}=1000\,\mathrm{N/m}$,
$k_{R,\mathrm{free}}=5\,\mathrm{Nm/rad}$,
$k_{R,\mathrm{stiff}}=200\,\mathrm{Nm/rad}$.
For this example the surface normal is \(n_s=[0,0,1]^\top\), with
\(t_1=[1,0,0]^\top\) and \(t_2=[0,1,0]^\top\). The surface-frame order follows
\cref{eq:surface_task_frame_matrix}, \(R_{\mathrm{surface}}=[\,t_1\;t_2\;n_s\,]\).
After transforming the tangent-first local gains to the base frame, the
base-frame translational stiffness is diagonal:

\begin{equation}
\label{eq:K_panda}
K_p =
\begin{bmatrix}
50   & 0  & 0    \\
0    & 50 & 0    \\
0    & 0  & 1000
\end{bmatrix}
\mathrm{N/m},
\qquad
K_R =
\begin{bmatrix}
5 & 0 & 0 \\
0 & 5 & 0 \\
0 & 0 & 200
\end{bmatrix}
\mathrm{Nm/rad}
\end{equation}

These are diagonal here \emph{only} because the surface normal and tangents
coincide with base-frame axes.
For a tilted wall with \(R_{\mathrm{surface}}\neq I\), applying
\(K_{p,0}=R_{\mathrm{surface}}K_{p,\mathrm{task}}R_{\mathrm{surface}}^\top\)
would produce a \textbf{full} $3\times3$ symmetric positive definite matrix
with non-zero
off-diagonal entries by the same congruence transformation used in
\cref{eq:general_task_frame_gain_transform}.

\subsection{\texorpdfstring{Damping Matrices $D_p$ and $D_R$}{Damping Matrices Dp and DR}}

For the numerical example, the damping matrices are selected as:

\begin{equation}
\label{eq:D_panda}
D_p =
\begin{bmatrix}
14.14 & 0     & 0     \\
0     & 14.14 & 0     \\
0     & 0     & 63.25
\end{bmatrix}
\mathrm{Ns/m}
\end{equation}

\begin{equation}
\label{eq:DR_panda}
D_R =
\begin{bmatrix}
4.47 & 0    & 0     \\
0    & 4.47 & 0     \\
0    & 0    & 28.28
\end{bmatrix}
\mathrm{Nm\,s/rad}
\end{equation}

Again, diagonal here only because the surface frame axes coincide with the
base-frame coordinate axes after the normal-first ordering is transformed.
For a general wall orientation all four matrices $K_p$,
$K_R$, $D_p$, $D_R$ would be full symmetric positive definite
$3\times3$ matrices in the base frame.

% ─────────────────────────────────────────────────────────────────
\section{Step 5 --- Pose Error and Impedance Wrench}
\label{app:panda_force}
% ─────────────────────────────────────────────────────────────────

\subsection{Position and Orientation Error}

Suppose the desired position and orientation are \(T_d\) with \(R_d=I_3\)
and $p_d=[0.136,\;0.545,\;0.395]^\top\,\mathrm{m}$
(the home position shifted $20\,\mathrm{mm}$ in $x$).
The current position and orientation at $\vq_0$ are \(T_{\mathrm{EE}}\) with \(R_{\mathrm{EE}}\) (from forward kinematics)
and $p_{\mathrm{EE}}=[0.116,\;0.545,\;0.395]^\top\,\mathrm{m}$.

Following \cref{eq:expanded_transformation_error}, the position-and-orientation error matrix is:

\begin{equation}
T_{\mathrm{err}} = T_d^{-1}T_{\mathrm{EE}} =
\begin{bmatrix}
R_d^\top R_{\mathrm{EE}} &
R_d^\top(p_{\mathrm{EE}} - p_d) \\
\mathbf{0}^\top & 1
\end{bmatrix}
=
\begin{bmatrix}
R_{\mathrm{EE}} & [-0.020,\;0,\;0]^\top \\
\mathbf{0}^\top & 1
\end{bmatrix}
\end{equation}

since \(R_d=I_3\) so \(R_d^\top=I_3\).

Reading off the sub-errors from the two blocks:

\paragraph{Relative translational displacement (upper-right column):}
\begin{equation}
\label{eq:panda_relative_displacement}
\Delta p_d = R_d^\top(p_{\mathrm{EE}} - p_d) =
\begin{bmatrix}
0.116 - 0.136 \\ 0 \\ 0
\end{bmatrix}
=
\begin{bmatrix}
-0.020 \\ 0 \\ 0
\end{bmatrix}
\mathrm{m}
\end{equation}

\paragraph{Positional control error:}
\begin{equation}
\label{eq:panda_error}
e_p = p_d-p_{\mathrm{EE}} =
\begin{bmatrix}
0.020 \\ 0 \\ 0
\end{bmatrix}
\mathrm{m}.
\end{equation}

\paragraph{Rotational error (upper-left block):}
\(\Delta R = R_d^\top R_{\mathrm{EE}} = R_{\mathrm{EE}}\).
Since the robot is at the desired orientation (\(R_{\mathrm{EE}}\approx R_d\)),
\(\Delta R\approx I\), giving
\(\phi = \arccos\!\bigl(\tfrac{\mathrm{tr}(I)-1}{2}\bigr) = 0\)
and $e_R = \mathbf{0}$.

\subsection{Impedance Wrench}

Assuming the robot is momentarily at rest
($\dot{p}_d=\mathbf{0}$, $\dot{p}_{\mathrm{EE}}=\mathbf{0}$, $\boldsymbol{\omega}_{\mathrm{EE}}=\mathbf{0}$),
the positional force and rotational moment from
\cref{eq:positional_impedance_force,eq:rotational_impedance_moment} are:

\begin{equation}
\label{eq:F_panda}
f = K_p\,e_p + D_p\left(\dot{p}_d-\dot{p}_{\mathrm{EE}}\right)
=
\begin{bmatrix}
50   & 0  & 0    \\
0    & 50 & 0    \\
0    & 0  & 1000
\end{bmatrix}
\begin{bmatrix}
0.020 \\ 0 \\ 0
\end{bmatrix}
=
\begin{bmatrix}
+1.0\,\mathrm{N} \\ 0 \\ 0
\end{bmatrix}
\end{equation}

\begin{equation}
m = K_R\,e_R = \mathbf{0}
\end{equation}

The combined forces and moments are:

\begin{equation}
F =
\begin{bmatrix}
f \\ m
\end{bmatrix}
=
\begin{bmatrix}
+1.0\,\mathrm{N} \\ 0 \\ 0 \\ 0 \\ 0 \\ 0
\end{bmatrix}
\in\mathbb{R}^6
\end{equation}

The \(+1\,\mathrm{N}\) in the \(x\)-direction is a restoring force pulling the
end-effector $20\,\mathrm{mm}$ back to the desired position.
The low in-plane stiffness $k_{\mathrm{free}}=50\,\mathrm{N/m}$ means this
restoring force is gentle: the wall yields compliantly under applied loads.

% ─────────────────────────────────────────────────────────────────
\section{Step 6 --- Joint Torques via \texorpdfstring{$\bm{\tau} = J^\top F$}{tau = J\^{}T F}}
\label{app:panda_torques}
% ─────────────────────────────────────────────────────────────────

Applying the control law \cref{eq:joint_torque_mapping} with $\mJ^\top \in \R{7\times6}$
from \cref{eq:J_panda}:

\begin{equation}
\label{eq:panda_torque_calc}
\bm{\tau} = \mJ^\top\,\vF
\end{equation}

Since $F = [1,\;0,\;0,\;0,\;0,\;0]^\top\,\mathrm{N}$,
only the first row of $\mJ$ (the $v_x$ row) contributes to $\vtau$:

\begin{equation}
\label{eq:panda_tau}
\bm{\tau} = \mJ^\top\,\vF =
\begin{bmatrix}
J_{11}\cdot 1 \\[2pt]
J_{21}\cdot 1 \\[2pt]
J_{31}\cdot 1 \\[2pt]
J_{41}\cdot 1 \\[2pt]
J_{51}\cdot 1 \\[2pt]
J_{61}\cdot 1 \\[2pt]
J_{71}\cdot 1
\end{bmatrix}
=
\begin{bmatrix}
(-0.5452)(1) \\[2pt]
(-0.5452)(1) \\[2pt]
( 0.0440)(1) \\[2pt]
(-0.3218)(1) \\[2pt]
( 0.0622)(1) \\[2pt]
( 0.0622)(1) \\[2pt]
( 0.0000)(1)
\end{bmatrix}
=
\begin{bmatrix}
-0.545\,\mathrm{Nm} \\[2pt]
-0.545\,\mathrm{Nm} \\[2pt]
+0.044\,\mathrm{Nm} \\[2pt]
-0.322\,\mathrm{Nm} \\[2pt]
+0.062\,\mathrm{Nm} \\[2pt]
+0.062\,\mathrm{Nm} \\[2pt]
\phantom{+}0.000\,\mathrm{Nm}
\end{bmatrix}
\end{equation}

\begin{itemize}
  \item $\tau_1 = \tau_2 = -0.545\,\mathrm{Nm}$:
        joints 1 and 2 contribute equally (consistent with their
        nearly identical Jacobian columns).
  \item $\tau_3 = +0.044\,\mathrm{Nm}$:
        joint 3 contributes a small opposing torque because its rotation
        axis ($z_3 \approx [0.707,\;0.707,\;0]^\top$) is
        at $45^\circ$ to the $x$-direction.
  \item $\tau_4 = -0.322\,\mathrm{Nm}$:
        the elbow joint has the largest
        translational lever arm to the end-effector.
  \item $\tau_5 = \tau_6 = +0.062\,\mathrm{Nm}$:
        the wrist joints contribute small opposing torques.
  \item $\tau_7 = 0$:
        joint 7's rotation axis ($-x_0$) is perpendicular to the
        applied force direction, so it cannot contribute to an $x$-directed
        Cartesian force.
\end{itemize}

% ─────────────────────────────────────────────────────────────────
\section{Step 7 --- Complete Control Loop Summary}
\label{app:panda_loop}
% ─────────────────────────────────────────────────────────────────

\begin{table}[htb]
\centering
\caption{Complete signal flow for one 1\,ms control cycle of the
         Panda impedance controller at the home configuration $\vq_0$.}
\label{tab:panda_loop}
\renewcommand{\arraystretch}{1.45}
\begin{tabular}{@{}c>{\raggedright\arraybackslash}p{3.6cm}>{\raggedright\arraybackslash}p{2.7cm}>{\raggedright\arraybackslash}p{5.9cm}@{}}
\toprule
Step & Quantity & Symbol & Value \\
\midrule
1  & Joint configuration   & $\vq$          & $[0, -45, 0, -135, 0, 90, 45]^\top$ deg \\
2  & EE position           & $p_{\mathrm{EE}}$ & $[0.116,\;0.545,\;0.395]^\top\,\mathrm{m}$ \\
3  & Jacobian ($6\times7$) & $J$   & See \cref{eq:J_panda} \\
4  & Error quantities & $\Delta p_d,e_p,e_R$ & $[-0.020,0,0]^\top$ m, $[0.020,0,0]^\top$ m, $\mathbf{0}$ \\
5  & Virtual inertia       & $M$   & $\mI_6$ \\
6  & Stiffness matrices    & $K_p, K_R$ & See \cref{eq:K_panda} \\
7  & Damping matrices      & $D_p, D_R$ & See \cref{eq:D_panda} \\
8  & Impedance forces and moments & $f,m$ & $[+1.0,0,0]^\top$ N, $\mathbf{0}$ \\
9  & Pseudo-inverse        & $J^+$ & See Eq.~(\ref{eq:Jp_panda}) \\
10 & Null-space projector  & $N_q$   & $\mI_7 - J^+\mJ$ \\
11 & Cartesian impedance torque & $\tau_{\mathrm{cart}}$ & $[-0.55, -0.55, +0.04, -0.32,$\newline$+0.06, +0.06, 0.00]^\top$ N\,m \\
12 & Null-space torque    & $\tau_{\mathrm{null}}$ & projected secondary torque; zero for this numerical example \\
\bottomrule
\end{tabular}
\end{table}

The C++ implementation of this control cycle is shown in Listing~\ref{lst:panda}:

\begin{lstlisting}[style=cppstyle,
  caption={One iteration of the Panda 7-DOF impedance control loop
           (\texttt{libfranka} torque callback, 1\,ms period)},
  label={lst:panda}]
// -- READ ROBOT STATE ---------------------------------------------
Eigen::Map<const Eigen::Matrix<double,7,1>> q(state.q.data());
Eigen::Map<const Eigen::Matrix<double,7,1>> dq(state.dq.data());

// -- STEP 2: End-effector position and orientation from libfranka --
Eigen::Affine3d T(Eigen::Matrix4d::Map(state.O_T_EE.data()));
Eigen::Vector3d p_EE  = T.translation();    // [3x1] position
Eigen::Matrix3d R_EE  = T.rotation();       // [3x3] orientation

// -- STEP 3: Jacobian J(q) from libfranka model --------------------
std::array<double,42> jac_arr =
    model.zeroJacobian(franka::Frame::kEndEffector, state);
Eigen::Map<const Eigen::Matrix<double,6,7>> J(jac_arr.data());

// -- STEP 4: Cartesian velocity  x_dot = J * q_dot ----------------
Eigen::Matrix<double,6,1> xdot = J * dq;

// (position-and-orientation error is computed in STEP 5 via T_err)

// -- STEP 5: Position and orientation error via T_err --------------
Eigen::Matrix3d R_d = T_des.rotation();
Eigen::Vector3d p_d = T_des.translation();

// Translational error in desired frame
Eigen::Vector3d e_p = p_d - p_EE;

// Rotational error: Delta_R = R_d^T * R_EE
Eigen::Matrix3d Delta_R = R_d.transpose() * R_EE;
double trace = Delta_R.trace();
double phi   = std::acos(std::clamp(0.5*(trace-1.0), -1.0, 1.0));
Eigen::Vector3d e_R = Eigen::Vector3d::Zero();
if (phi > 1e-6)
    e_R = (phi/(2.0*std::sin(phi))) *
          Eigen::Vector3d(Delta_R(2,1)-Delta_R(1,2),
                          Delta_R(0,2)-Delta_R(2,0),
                          Delta_R(1,0)-Delta_R(0,1));

// -- STEPS 6-8: Full SPD impedance matrices -----------------------
// Surface-frame order: [normal, tangent_1, tangent_2].
// The transformed base-frame matrices are diagonal here only because
// the surface axes coincide with base-frame coordinate axes.
Eigen::Matrix3d Kp = Eigen::Vector3d{50,50,1000}.asDiagonal();
Eigen::Matrix3d KR = Eigen::Vector3d{5,5,200}.asDiagonal();
Eigen::Matrix3d Dp = Eigen::Vector3d{14.14,14.14,63.25}.asDiagonal();
Eigen::Matrix3d DR = Eigen::Vector3d{4.47,4.47,28.28}.asDiagonal();
// For a tilted wall:
// Kp = R_surface * Kp_task * R_surface.transpose() (full matrix)

// Cartesian velocity split
Eigen::Vector3d pdot  = xdot.head<3>();
Eigen::Vector3d pdot_d = Eigen::Vector3d::Zero();
Eigen::Vector3d omega = xdot.tail<3>();

// Positional force and rotational moment
Eigen::Vector3d f = Kp * e_p + Dp * (pdot_d - pdot);
Eigen::Vector3d m = KR * e_R - DR * omega;

// Combined 6D forces and moments
Eigen::Matrix<double,6,1> F;
F.head<3>() = f;  F.tail<3>() = m;

// -- STEP 9-10: Pseudo-inverse and null-space ---------------------
// Pseudo-inverse and null-space projector.
// The primary control law needs only J.transpose(); N is used for
// the projected secondary null-space torque.
Eigen::Matrix<double,7,6> Jp =
    J.transpose() * (J * J.transpose()).inverse();    // J+ [7x6]
Eigen::Matrix<double,7,7> N =
    Eigen::Matrix<double,7,7>::Identity() - Jp * J;

// -- STEP 11: Primary torque  tau = J^T * F -----------------------
//  (J^T is exact, no pseudo-inverse needed here)
Eigen::Matrix<double,7,1> tau_task = J.transpose() * F;

// -- NULL-SPACE TORQUE --------------------------------------------
// In the implemented controller this term is computed from the
// projected null-space objective. It is zero in this numerical example.
Eigen::Matrix<double,7,1> tau_null = Eigen::Matrix<double,7,1>::Zero();

// -- GRAVITY + CORIOLIS COMPENSATION ------------------------------
Eigen::Map<const Eigen::Matrix<double,7,1>> tau_g(state.gravity.data());
Eigen::Map<const Eigen::Matrix<double,7,1>> tau_c(state.coriolis.data());

// -- FINAL COMMAND ------------------------------------------------
// Final command: primary impedance + projected null-space torque
// + gravity + Coriolis compensation
Eigen::Matrix<double,7,1> tau_cmd =
    tau_task + tau_null + tau_g + tau_c;

// -- SAFETY CLIPPING ----------------------------------------------
const std::array<double,7> tau_max = {87,87,87,87,12,12,12};
std::array<double,7> tau_out{};
for (size_t i=0; i<7; ++i)
  tau_out[i] = std::max(-tau_max[i], std::min(tau_max[i], tau_cmd(i)));

return franka::Torques(tau_out);
\end{lstlisting}

% ─────────────────────────────────────────────────────────────────
\section{Step 8 --- Effect of Stiffness on Torque Distribution}
\label{app:panda_stiffness}
% ─────────────────────────────────────────────────────────────────

To illustrate how $K_p$ scales the torque command, the
normal stiffness $k_{\mathrm{stiff}}$ is varied while keeping the
same $x$-direction error of $20\,\mathrm{mm}$
(this scalar example varies the stiffness in the displayed force direction):

\begin{center}
\renewcommand{\arraystretch}{1.35}
\begin{tabular}{@{}lrrrrrrr@{}}
\toprule
$k_{\mathrm{stiff}}\;[\mathrm{N/m}]$
  & $\tau_1$ & $\tau_2$ & $\tau_3$ & $\tau_4$
  & $\tau_5$ & $\tau_6$ & $\tau_7$ [Nm] \\
\midrule
$200$
  & $+0.218$ & $+0.218$ & $-0.018$ & $+0.129$
  & $-0.025$ & $-0.025$ & $0.000$ \\
$500$
  & $+0.545$ & $+0.545$ & $-0.044$ & $+0.322$
  & $-0.062$ & $-0.062$ & $0.000$ \\
$1000$
  & $+1.090$ & $+1.090$ & $-0.088$ & $+0.644$
  & $-0.124$ & $-0.124$ & $0.000$ \\
\bottomrule
\end{tabular}
\end{center}

\paragraph{Key observations.}
\begin{itemize}
  \item The torque magnitudes scale linearly with $K_x$, since
        $F_x = K_x \times 0.020$ and $\tau = \mJ^\top\vF$.
  \item The \emph{ratio} between the torques of different joints is
        constant across all three rows: it depends only on the Jacobian
        $\mJ$ at the current configuration, not on $\mK$.
  \item Joint\,7 always produces zero torque for an $x$-directed force
        at this configuration, regardless of stiffness.
\end{itemize}

\bigskip
\noindent
This completes the 7-DOF Panda worked example.
Every quantity computed here—from the DH chain to the final torque vector—
corresponds directly to one iteration of the C++ callback in
Listing~\ref{lst:panda} and is executed 1000 times per second during
operation.


%% =============================================================
%%  APPENDIX B – DATA LOGGING AND ANALYSIS WORKFLOW
%% =============================================================
